\section{Experimental evaluation}
\label{sec:eval}

The bounds proved so far are worst-case statements, and a worst-case statement
says nothing about how often the worst case arrives. \Cref{thm:deletebound}
permits a deletion to rotate at every level of the tree, and \Cref{fig:cascade}
exhibits a tree where three levels do rotate, but neither result indicates
whether a program that deletes a million keys should expect to pay for that
behaviour. This section measures what actually happens.

\subsection{Method}

We implemented the AVL tree of \Cref{alg:rotate,alg:rebalance,alg:insert,alg:delete}
and, for comparison, an ordinary binary search tree with no balancing. The two
implementations share their node representation and their search routine, so any
difference between them comes from the balancing and not from unrelated
bookkeeping.

We measure operation counts rather than elapsed time. A comparison count and a
rotation count are properties of the structure, while a wall-clock reading also
records the memory hierarchy, the allocator, and the language runtime, none of
which the analysis of the preceding sections speaks about. Counting comparisons
also makes the results reproducible on any machine.

Three quantities are recorded. The first is the height of the tree after all keys
have been inserted. The second is the mean number of comparisons for a successful
search, which equals the mean node depth plus one and which we obtain by summing
the depths of all nodes. The third is the rebalancing work per update, counted in
two ways: rebalancing \emph{events}, where a double rotation counts once, and
single rotations, where a double rotation counts twice. The first count measures
how often the repair runs, and the second measures how much pointer work it does.

Each tree is built by inserting $n$ distinct keys, and for the deletion
measurements every key is then removed in a fresh random order until the tree is
empty. Two key orders are used: a uniformly random permutation, and increasing
order, which is the adversarial case of \Cref{fig:degenerate}. Sizes run from
$100$ to $10^{6}$ keys. Results are averaged over twenty independent trials for
$n \le 3000$, five trials up to $10^{5}$, and a single trial at the largest
sizes, where the variation between trials is negligible.

\subsection{Height against the bound}

\Cref{fig:plotheight} plots the measured height against the two bounds of
\Cref{thm:heightbound}. The horizontal axis is logarithmic, so all three series
are straight lines and their slopes are directly comparable.

\begin{figure}[t]
  \centering
  % Measured height against the two analytic bounds. The horizontal axis is
% logarithmic, which is why all three series come out as straight lines.
%
% The gap that matters is the one between the measured series and the upper
% bound: the bound is correct but it is not tight, because the minimum-node
% trees that attain it are vanishingly unlikely to arise from random keys.
\begin{tikzpicture}
  \begin{axis}[
      avlaxis,
      xmode        = log,
      log basis x  = 10,
      xlabel       = {number of keys $n$},
      ylabel       = {height $h$},
      xmin         = 8, xmax = 1.4e6,
      ymin         = 0, ymax = 30,
      legend pos   = north west,
    ]

    \addplot[avlbound] table[x = n, y = upper] {data/height_vs_n.dat};
    \addlegendentry{upper bound $1.4404\lgg(n+2)-1.3277$}

    \addplot[avlmeasured] table[x = n, y = meanh] {data/height_vs_n.dat};
    \addlegendentry{measured height}

    \addplot[avlfloor] table[x = n, y = lower] {data/height_vs_n.dat};
    \addlegendentry{lower bound $\lgg(n+1)$}

  \end{axis}
\end{tikzpicture}

  \caption{Height of an AVL tree built from random keys, against the analytic
  bounds. The measured series runs close to the information-theoretic lower
  bound rather than to the upper bound.}
  \label{fig:plotheight}
\end{figure}

The measured height sits much nearer the lower bound than the upper one. At
$n = 10^{6}$ the tree has height $23$, against a lower bound of $19.93$ and an
upper bound of $27.38$. Trees grown from random keys use about $15$ percent more
levels than a perfect tree, not the $44$ percent the worst case permits.

The gap does not indicate that \Cref{thm:heightbound} is loose. The bound is
attained only by the minimum-size trees, and those are rare among the trees a
random key order produces: of the $54$ nodes in \Cref{fig:cascade}, every single
one has a balance factor of $\pm 1$, and a random sequence of insertions has
essentially no chance of assembling such a tree. The bound describes an adversary,
and the measurements describe an average.

One further observation concerns variance. At every size from $100$ keys upward,
every trial produced trees of identical height. The height of an AVL tree is
determined by its size much more tightly than the two-bound gap suggests.

\subsection{Search cost against an unbalanced tree}

\Cref{fig:plotsearch} compares the AVL tree with the unbalanced tree on both key
orders. Both axes are logarithmic, because the four series do not share an order
of growth.

\begin{figure}[t]
  \centering
  % Average comparisons for a successful search. Both axes are logarithmic,
% because the four series do not share an order of growth: three of them are
% logarithmic and one is linear.
%
% The separation is the whole argument for balancing. On sorted keys the plain
% tree needs hundreds of times more comparisons than the balanced one, and the
% factor grows with every key added.
\begin{tikzpicture}
  \begin{axis}[
      avlaxis,
      xmode        = log,
      ymode        = log,
      log basis x  = 10,
      xlabel       = {number of keys $n$},
      ylabel       = {mean comparisons per search},
      xmin         = 80, xmax = 1.3e4,
      ymin         = 4, ymax = 1.2e4,
      % Below the axis: every corner inside the frame has a curve running
      % through it, so an inset legend would sit on top of the data.
      legend style = {at = {(0.5, -0.30)}, anchor = north, draw = rfaint},
      legend columns = 2,
      legend cell align = left,
    ]

    \addplot[mark = *, mark size = 1.5pt, thick, color = rmid, dashed]
      table[x = n, y = bstsorted] {data/search_cost.dat};
    \addlegendentry{plain tree, sorted keys}

    \addplot[avlsecond] table[x = n, y = bstrand] {data/search_cost.dat};
    \addlegendentry{plain tree, random keys}

    \addplot[avlmeasured] table[x = n, y = avlrand] {data/search_cost.dat};
    \addlegendentry{AVL tree, random keys}

    \addplot[no marks, thick, dotted, color = rkink]
      table[x = n, y = avlsorted] {data/search_cost.dat};
    \addlegendentry{AVL tree, sorted keys}

  \end{axis}
\end{tikzpicture}

  \caption{Mean comparisons for a successful search. Three series are
  logarithmic in $n$ and one is linear. The unbalanced tree on sorted keys is the
  failure case the AVL condition exists to prevent.}
  \label{fig:plotsearch}
\end{figure}

On random keys the two structures are close. At $n = 10^{4}$ the AVL tree needs
$12.59$ comparisons and the unbalanced tree needs $16.98$, a penalty of about a
third. A random binary search tree has expected height $\Theta(\log n)$ already,
so balancing buys a constant factor and no more.

On sorted keys the comparison is different in kind. The AVL tree needs $12.36$
comparisons, marginally fewer than on random keys, because inserting in order
produces a tree that is nearly perfectly balanced. The unbalanced tree needs
$5000.5$ comparisons, which is $(n+1)/2$ and confirms that it has degenerated
into the list of \Cref{fig:degenerate}. The balanced structure is faster by a
factor of $404$ at ten thousand keys, and the factor grows linearly with $n$.

The practical reading is that balancing is not an optimisation of the average
case. It is insurance against an input order that costs nothing to produce and
that ordinary workloads produce by accident.

\subsection{Rebalancing work per operation}

\Cref{tab:rotations} reports the repair work, and \Cref{fig:plotrotations} plots
the same figures.

\begin{table}[t]
  \centering
  \caption{Rebalancing work per update on random keys. A rebalancing event is one
  repair, counting a double rotation once; the rotation columns count a double
  rotation as two. All four quantities are flat in $n$.}
  \label{tab:rotations}
  \begin{tabular}{@{}S[table-format=7.0]S[table-format=1.4]S[table-format=1.4]S[table-format=1.4]S[table-format=1.4]@{}}
    \toprule
    & \multicolumn{2}{c}{Per insertion} & \multicolumn{2}{c}{Per deletion} \\
    \cmidrule(lr){2-3}\cmidrule(lr){4-5}
    {$n$} & {events} & {rotations} & {events} & {rotations} \\
    \midrule
    100    & 0.4510 & 0.6840 & 0.2360 & 0.3360 \\
    300    & 0.4623 & 0.6843 & 0.2600 & 0.3623 \\
    1000   & 0.4675 & 0.7011 & 0.2587 & 0.3618 \\
    3000   & 0.4595 & 0.6893 & 0.2636 & 0.3714 \\
    10000  & 0.4658 & 0.6981 & 0.2642 & 0.3691 \\
    30000  & 0.4629 & 0.6946 & 0.2667 & 0.3756 \\
    100000 & 0.4668 & 0.6993 & 0.2661 & 0.3727 \\
    300000 & 0.4642 & 0.6963 & 0.2661 & 0.3732 \\
    \bottomrule
  \end{tabular}
\end{table}

\begin{figure}[t]
  \centering
  % Rebalancing work per operation, against a logarithmic size axis.
%
% Both series are flat. That is the point of the picture: the analysis of
% Section 6 allows the number of rotations per deletion to grow with the
% height, and measurement says the average does not. The worst case and the
% average case part company here.
\begin{tikzpicture}
  \begin{axis}[
      avlaxis,
      height       = 5.6cm,
      xmode        = log,
      log basis x  = 10,
      xlabel       = {number of keys $n$},
      ylabel       = {per operation},
      xmin         = 80, xmax = 4e5,
      ymin         = 0, ymax = 1.25,   % headroom, so the legend clears the data
      legend pos   = north east,
      legend columns = 2,
      legend cell align = left,
    ]

    \addplot[avlmeasured] table[x = n, y = insrot] {data/rotations.dat};
    \addlegendentry{single rotations, insertion}

    \addplot[avlsecond] table[x = n, y = delrot] {data/rotations.dat};
    \addlegendentry{single rotations, deletion}

    \addplot[no marks, thick, dashed, color = rkink]
      table[x = n, y = insev] {data/rotations.dat};
    \addlegendentry{rebalance events, insertion}

    \addplot[no marks, thick, dotted, color = rmid]
      table[x = n, y = delev] {data/rotations.dat};
    \addlegendentry{rebalance events, deletion}

  \end{axis}
\end{tikzpicture}

  \caption{Rebalancing work per update against tree size. Every series is flat
  across three and a half orders of magnitude, including the deletion series that
  \Cref{thm:deletebound} allows to grow with the height.}
  \label{fig:plotrotations}
\end{figure}

Two results stand out. The first is that insertion averages $0.464$ rebalancing
events, so roughly one insertion in two triggers no repair at all. This is
consistent with \Cref{thm:insertone}, which caps the count at one, and it
locates the average at about half the cap.

The second result is the one \Cref{thm:deletebound} does not predict. Deletion
averages $0.266$ rebalancing events, fewer than insertion rather than more, and
the figure does not grow as $n$ increases from $100$ to $300\,000$. The worst
case of $h$ repairs per deletion is real, as \Cref{fig:cascade} demonstrates, but
it requires a tree with no slack at any level, and ordinary trees have slack
nearly everywhere. A cascade stops as soon as it reaches a node whose balance
factor was $0$, and after a repair that node absorbs the lost level without
passing it on. In a random tree such a node is never far away.

Deletion does less repair work than insertion for a related reason. Most nodes in
a tree of any size are near the leaves, so a randomly chosen key is usually
shallow in the sense that matters here, and removing it disturbs few ancestors.

\subsection{Threats to validity}

Three limitations qualify the results above.

The measurements count comparisons and rotations, not time. A rotation writes
three pointers and may evict cache lines, so a structure that rotates less can be
faster in practice even when it compares more. \Cref{sec:comparison} returns to
this point with timing results from the literature.

The key orders tested are a uniformly random permutation and increasing order.
Real workloads sit between the two, and a partially sorted input, such as records
that are mostly ordered with occasional insertions out of place, is not covered
here. The unbalanced tree degrades smoothly between the two series of
\Cref{fig:plotsearch}, so its behaviour on such input lies between them, but we
did not measure where.

The deletion measurements remove every key from the tree. A workload that
interleaves insertions and deletions at a steady size may reach a different
steady state, since the trees it produces are shaped by both operations rather
than built by one and dismantled by the other.
